\usepackage{xeCJK}
\usepackage{geometry}
\geometry{left=2cm,right=2cm,top=2.5cm,bottom=2.5cm}
\setlength{\headheight}{12.7pt}

\usepackage{titlesec}
\titleformat{\section}
  {\normalfont\large\bfseries}
  {\thesection}
  {1em}
  {}
\titlespacing*{\section}
  {0pt}
  {3.5ex plus 1ex minus .2ex}
  {2.3ex plus .2ex}

\usepackage{amsmath}
\usepackage{amssymb}
\usepackage{amsthm}
\usepackage{enumitem}
\setlist[enumerate]{itemsep=0pt,topsep=0pt,parsep=0pt,leftmargin=0pt,itemindent=2.5em}
\setlist[itemize]{itemsep=0pt,topsep=0pt,parsep=0pt,leftmargin=0pt,itemindent=2.5em}
\usepackage{fancyhdr}
\pagestyle{fancy}
\fancyhf{}
\usepackage{graphicx}
\usepackage{hyperref}
\usepackage{indentfirst}
\usepackage{lmodern}


\begin{document}
\setlength{\abovedisplayskip}{1pt}
\setlength{\belowdisplayskip}{1pt}

\section{Łoś-Vaught 判别法}

\hypertarget{sec:定义1.1}{}
\noindent\textbf{定义1.1 ($\kappa$-定言)}

给定一阶语言$\mathcal{L}$和$\mathcal{L}$-理论$T$.
如果$T$的任意两个势为$\kappa$的模型都同构, 那么称$T$是$\kappa$-\textbf{定言}的.

\vspace{10pt}

\hypertarget{sec:定理1.1}{}
\noindent\textbf{定理1.1 (Łoś-Vaught判别法)}

给定一阶语言$\mathcal{L}$和$\mathcal{L}$-理论$T$.

如果$T$的模型都是无限的,
并且在某个基数$\kappa\geqslant|\mathcal{L}|$上是$\kappa$-定言的, 那么$T$
是\href{../02 一阶语义/03 一阶理论与模型#nameddest=sec:定义2.1}{完备}的.

\noindent{\kaishu 证明.}

任取$T$的模型$A$, 依假设$A$是无限的.
\begin{enumerate}
    \item 如果$|A|<\kappa$,
    那么根据\href{03 Löwenheim-Skolem 定理#nameddest=sec:定理2.1}
    {向上Löwenheim-Skolem定理}, 可以取一个结构$A'$使得
    \[
        A\preccurlyeq A'\quad\land\quad|A'|=\kappa,
    \]
    \item 如果$|A|=\kappa$, 那么令$A'=A$,
    \item 如果$|A|>\kappa$,
    那么根据\href{03 Löwenheim-Skolem 定理#nameddest=sec:定理1.1}
    {向下Löwenheim-Skolem定理}, 可以取一个结构$A'$使得
    \[
        A'\preccurlyeq A\quad\land\quad|A'|=\kappa.
    \]
\end{enumerate}

综上, 存在结构$A'$使得$A\equiv A'$并且$|A'|=\kappa$. 显然$A'$也是$T$的模型.

同理, 任取$T$的模型$B$, 也存在模型$B'$使得$B\equiv B'$并且$|B'|=\kappa$.
依假设可知$A'\cong B'$, 于是
\[
    A\equiv A'\equiv B'\equiv B.
\]

% 这里可以看出, 判定的前提只需要势为$\kappa$的模型都初等等价, 不需要同构.
% 这样写应该是因为很多时候同构虽然更强, 但是更好判定. 于是应用定理比较方便.

所以$T$的模型都初等等价,
由\href{../02 一阶语义/03 一阶理论与模型#nameddest=sec:定理2.1}{引理}即得
$T$是完备的. \hfill $\qedsymbol$

\end{document}